% ==============================================================================
% KATA PENGANTAR
% Sesuai Format Pedoman Skripsi FST UNIBA Madura 2026
% ==============================================================================

\clearpage
\phantomsection
\pagestyle{frontstyle}
\addcontentsline{toc}{chapter}{KATA PENGANTAR}

\begin{center}
    {\fontsize{14pt}{18pt}\selectfont \bfseries KATA PENGANTAR}
\end{center}

\vspace{0.3cm}

\textit{Alhamdulillahi Rabbil ’Alamin}, segala puja dan puji syukur penulis panjatkan kehadirat Allah SWT atas berkah, rahmat, dan taufik-Nya, sehingga penulis dapat menyelesaikan penyusunan proposal skripsi yang berjudul \textbf{“\skripsiJudul”} ini dengan baik. Naskah proposal ini disusun sebagai salah satu syarat untuk memenuhi ketentuan dalam kurikulum serta menyelesaikan jenjang studi guna memperoleh gelar \skripsiGelar\ (\skripsiSingkatanGelar) pada Program Studi \skripsiProdi, \skripsiFakultas, \skripsiUniversitas.

Keberhasilan dalam merumuskan dan menyelesaikan proposal skripsi ini tentunya tidak terlepas dari bimbingan, petunjuk, arahan, saran, serta dukungan moral yang sangat berharga dari berbagai pihak. Oleh karena itu, melalui lembar ini, penulis menyampaikan rasa hormat dan terima kasih yang setulus-tulusnya kepada:

\begin{enumerate}\singlespacing\setlength{\itemsep}{1.5pt}\setlength{\parsep}{0pt}\setlength{\topsep}{2pt}
    \item Bapak \skripsiDekan, selaku Dekan Fakultas Sains dan Teknologi Universitas K.H. Bahaudin Mudhary Madura.
    \item Bapak Mohammad Iqbal Bachtiar, S.T., M.T., selaku Wakil Dekan Fakultas Sains dan Teknologi Universitas K.H. Bahaudin Mudhary Madura.
    \item Bapak \skripsiKaprodi, selaku Ketua Program Studi \skripsiProdi\ Universitas K.H. Bahaudin Mudhary Madura atas bimbingan dan dorongan akademis yang senantiasa diberikan.
    \item Bapak \skripsiPembimbingSatu, selaku Dosen Pembimbing I yang telah meluangkan banyak waktu, tenaga, pemikiran, serta arahan ilmiah yang kritis dan konstruktif dalam penyempurnaan gagasan penelitian ini.
    \item Bapak \skripsiPembimbingDua, selaku Dosen Pembimbing II atas ketelitian, bimbingan teknis perancangan metodologi, saran metodis, dan motivasi yang senantiasa diberikan kepada penulis.
    \item Seluruh Bapak dan Ibu Dosen serta Tenaga Kependidikan Program Studi \skripsiProdi\ Fakultas Sains dan Teknologi Universitas K.H. Bahaudin Mudhary Madura atas transfer ilmu pengetahuan dan fasilitas akademik selama perkuliahan.
    \item Teristimewa kepada Kedua Orang Tua tercinta, Ayahanda dan Ibunda, serta segenap keluarga besar atas doa yang tiada henti, cinta kasih, keikhlasan, dan pengorbanan tanpa batas yang menjadi sumber kekuatan utama bagi penulis.
    \item Rekan-rekan mahasiswa seperjuangan Program Studi \skripsiProdi\ angkatan 2023 atas diskusi, kebersamaan, dan dukungan semangat selama proses perkuliahan dan riset.
    \item Seluruh pihak yang telah memberikan kontribusi, baik secara langsung maupun tidak langsung, yang tidak dapat penulis sebutkan satu per satu.
\end{enumerate}

Penulis menyadari bahwa naskah proposal skripsi ini masih memiliki keterbatasan dan ketidaksempurnaan. Oleh sebab itu, penulis senantiasa terbuka menerima kritik objektif, evaluasi ilmiah, serta saran konstruktif dari dosen penguji dan pembaca demi penyempurnaan naskah dan pelaksanaan eksperimen penelitian ke depan. Semoga karya ini mampu memberikan kontribusi nyata bagi pengembangan ilmu pengetahuan di bidang pemrosesan bahasa alami daerah berdaya komputasi rendah (\textit{low-resource language processing}) serta bermanfaat bagi masyarakat.

\vspace{0.3cm}

\begin{flushright}
\begin{tabular}{l}
    \skripsiKota, \dots\ \dots\ \skripsiTahun\\[1.4cm]
    \skripsiPenulis
\end{tabular}
\end{flushright}
